% !TeX program = xelatex
% ============================================================
% pmi/pmi.tex — Программа и методика испытаний
% ГОСТ 19.301-79 «Программа и методика испытаний.
% Требования к содержанию и оформлению»
% ============================================================
\documentclass[12pt,a4paper]{extarticle}

\input{../common/preamble}

\newcommand{\docname}{Test Program and Methodology}
\newcommand{\doccode}{\hseDocCodeBase\ 51 01-1}
\newcommand{\doccodelu}{\hseDocCodeBase\ 51 01-1-ЛУ}

\input{../common/commands}
\input{../common/styles}

\begin{document}
\sloppy

% ── Титульный лист ──────────────────────────────────────────
\input{../common/title_template}


% ============================================================
%  ENGLISH DOCUMENT BODY (project.lang == en)
%  ГОСТ 19.301-79 Test Program and Methodology — translated per
%  §2.3.4.2 / §2.5.3.1. Structure mirrors the Russian version 1:1;
%  references follow IEEE style. Fill the yellow placeholders.
% ============================================================

% ============================================================
% ANNOTATION
% ============================================================
\section*{ANNOTATION}
\addcontentsline{toc}{section}{ANNOTATION}

This Test Program and Methodology for the development of the program \enquote{\hseProjectName} contains the following sections: \enquote{Test Object}, \enquote{Test Objective}, \enquote{Requirements for the Program}, \enquote{Requirements for the Program Documentation}, \enquote{Test Facilities and Procedure}, \enquote{Test Methods}.

The \enquote{Test Object} section specifies the name, field of application and designation of the program under test.

The \enquote{Test Objective} section states the purpose of the testing.

The \enquote{Requirements for the Program} section contains the list of requirements from the Statement of Work that are subject to verification during testing.

The \enquote{Requirements for the Program Documentation} section lists the program documentation submitted for testing.

The \enquote{Test Facilities and Procedure} section contains information on the hardware and software used during testing, as well as the procedure for their use.

The \enquote{Test Methods} section provides a description of the methods used and the specific checks.

% Подсказка: перечислите, что приведено в приложениях вашей ПМИ.
The appendices provide the list of terms, the traceability matrix of requirements to automated tests and the list of supporting test materials.

\clearpage

% ============================================================
% LIST OF APPLICABLE STANDARDS
% ============================================================
\noindent
This document has been developed in accordance with the requirements of:

\begin{enumerate}[label=\arabic*), leftmargin=1.25cm]
\item GOST 19.101-77 Types of programs and program documents\cite{gost19101}.
\item GOST 19.103-77 Designation of programs and program documents\cite{gost19103}.
\item GOST 19.104-78 Basic inscriptions\cite{gost19104}.
\item GOST 19.105-78 General requirements for program documents\cite{gost19105}.
\item GOST 19.106-78 Requirements for program documents produced
      in printed form\cite{gost19106}.
\item GOST 19.301-79 Test program and methodology. Requirements
      for content and presentation\cite{gost19301}.
\end{enumerate}

Changes to this document are made in accordance with GOST\,19.604-78\cite{gost19604}.

\clearpage

% ============================================================
% CONTENTS
% ============================================================
\hypersetup{linkcolor=black}
\tableofcontents
\hypersetup{linkcolor=blue}
\thispagestyle{fancy}
\clearpage

% ============================================================
% 1. TEST OBJECT
% ============================================================
\section{TEST OBJECT}

\subsection{Name of the program under test}

The name of the program~--- \enquote{\projectname}.

% Английское наименование подставляется из настроек проекта (поле
% «Название работы на английском»).
The name of the program in English~--- \enquote{\hseProjectNameEn}.

\subsection{Field of application of the program under test}

% Подсказка: кратко опишите назначение и область применения программы,
% а также какой контур системы рассматривается в рамках испытаний.

The program \enquote{\projectname} is \hseFill{brief description of the system purpose}. The system covers \hseFill{main domains and subsystems}.

Within the scope of this document, \hseFill{the tested contour of the system} is considered.

\subsection{Designation of the program under test}

The name of the development topic~--- \enquote{\projectname}.

% Краткое наименование подставляется из настроек проекта (поле
% «Краткое наименование программы»).
The short name of the program~--- \enquote{\hseProjectNameShort}.

\clearpage

% ============================================================
% 2. TEST OBJECTIVE
% ============================================================
\section{TEST OBJECTIVE}

The objective of the testing~--- to verify the conformance of the characteristics of the developed program with the functional and other requirements set out in the program documentation \enquote{Statement of Work}.

\clearpage

% ============================================================
% 3. REQUIREMENTS FOR THE PROGRAM
% ============================================================
\section{REQUIREMENTS FOR THE PROGRAM}

\subsection{Requirements for functional characteristics}

\subsubsection{Composition of the functions performed}

% Подсказка: идентификаторы (ТЗ-Ф-01 и т.д.) связывают требования
% с тестами в панели «Трассировка требований».

The functional characteristics of the program listed in the requirements table of the Statement of Work are subject to verification. Each functional requirement with an identifier of the form \texttt{ТЗ-Ф-XX} is matched with an automated test, a test suite or another check.

\input{../technical_specification/requirements_table}

\clearpage

\subsubsection{Organization of input data}

% Подсказка: перечислите проверяемые входные данные и форматы
% вашей системы (структура запросов, валидация, ограничения).
\begin{hseExample}
During testing, compliance with the Statement of Work requirements for the organization of input data is verified. The following are subject to verification:

\begin{itemize}[leftmargin=1.25cm]
    \item acceptance of \hseFill{type of requests or input data} with data transmission in the \hseFill{input data format} format;
    \item validation of incoming data against the schemas and data types specified in the interface specification;
    \item compliance with the size limits \hseFill{type of input data limits};
    \item \hseFill{other input data organization checks}.
\end{itemize}

Verification is performed both on valid input data and on erroneous data sets intended to confirm the operation of validation.
\end{hseExample}

\subsubsection{Organization of output data}

% Подсказка: перечислите проверяемые выходные данные (коды ответов,
% форматы, структура ошибок).
\begin{hseExample}
During testing, compliance with the Statement of Work requirements for the organization of output data is verified. The following are subject to verification:

\begin{itemize}[leftmargin=1.25cm]
    \item return of results in the \hseFill{output data format} format;
    \item use of \hseFill{status codes or response statuses} for successful and erroneous responses;
    \item return of a structured error description sufficient for diagnostics without disclosing confidential data;
    \item \hseFill{other output data organization checks}.
\end{itemize}
\end{hseExample}

\subsubsection{Requirements for timing characteristics}

% Подсказка: перечислите временные показатели из ТЗ, подлежащие
% проверке (время отклика, пропускная способность).
\begin{hseExample}
During testing, the timing characteristics established in the Statement of Work are verified:

\begin{itemize}[leftmargin=1.25cm]
    \item the processing time of \hseFill{type of requests or operations} must not exceed \hseFill{target response time};
    \item the specified indicators must be maintained under normal load of up to~\hseFill{target number of users};
    \item the target request rate must correspond to the range \hseFill{target request rate}.
\end{itemize}
\end{hseExample}

\subsection{Requirements for the interface}

% Подсказка: опишите проверяемые требования к программному интерфейсу
% (методы API, контракты, спецификации).
\begin{hseExample}
During testing, compliance with the Statement of Work requirements for the program interface is verified. The following are subject to verification:

\begin{itemize}[leftmargin=1.25cm]
    \item implementation of \hseFill{type of program interface} in accordance with \hseFill{interface principles and specifications};
    \item availability of the current interface specification \hseFill{specification location};
    \item verification of access to protected functions based on \hseFill{authentication and authorization mechanism};
    \item uniformity of the response structure, data formats and field naming conventions.
\end{itemize}
\end{hseExample}

\subsection{Requirements for marking and versioning}

% Подсказка: скорректируйте список под требования вашего ТЗ,
% если они заданы.
\begin{hseExample}
During testing, compliance with the Statement of Work requirements for the marking and versioning of system components is verified. The following are subject to verification:

\begin{itemize}[leftmargin=1.25cm]
    \item conformance of program code versions to the \hseFill{versioning scheme used} specification;
    \item presence of version tags \hseFill{type of build artifacts};
    \item presence of build metadata with the build date and the source code version identifier;
    \item storage of source code, configurations and build artifacts in repositories that support version control.
\end{itemize}
\end{hseExample}

\clearpage

% ============================================================
% 4. REQUIREMENTS FOR THE PROGRAM DOCUMENTATION
% ============================================================
\section{REQUIREMENTS FOR THE PROGRAM DOCUMENTATION}

\subsection{Composition of the program documentation submitted for testing}

% Подсказка: перечислите документы, предъявляемые на испытания,
% с указанием соответствующих ГОСТ.

The following program documentation must be submitted for testing:
\begin{enumerate}[label=\arabic*), leftmargin=1.25cm]
    \item \enquote{\projectname}. Statement of Work (GOST 19.201-78).
    \item \enquote{\projectname}. Source Listing (GOST 19.401-78).
    \item \enquote{\projectname}. Programmer's Manual (GOST 19.504-79).
    \item \enquote{\projectname}. Test Program and Methodology
          (GOST 19.301-79).
\end{enumerate}

\subsection{Special requirements}

The program documentation must be produced in accordance with GOST 19.106-78\cite{gost19106} and the GOST standards for each type of document (see~clause~4.1).

% Подсказка: при необходимости укажите иные специальные требования
% к оформлению документации.
Supporting test materials may be formatted as appendices to this document.

\clearpage

% ============================================================
% 5. TEST FACILITIES AND PROCEDURE
% ============================================================
\section{TEST FACILITIES AND PROCEDURE}

\subsection{Hardware used during testing}

% Подсказка: опишите аппаратное окружение испытаний (ОС, процессор,
% память) и контуры запуска.
\begin{hseExample}
Testing is carried out on hardware with the following characteristics: \hseFill{OS, processor, memory}. The following test environments are used:

% Таблица набрана через \captionof (не \begin{table}[H]): плавающие
% окружения внутри жёлтого блока-образца недопустимы.
\begin{center}
\small
\captionof{table}{Test environments}
\label{tab:pmi_test_environments}
\begin{tabularx}{\textwidth}{|>{\raggedright\arraybackslash}p{0.21\textwidth}|>{\raggedright\arraybackslash}p{0.29\textwidth}|X|}
\hline
\textbf{Environment} & \textbf{Composition} & \textbf{Purpose} \\
\hline
Local environment & \hseFill{composition of the local environment} & Quick test runs without full deployment of the environment. \\
\hline
\hseFill{integration environment} & \hseFill{composition of the integration environment} & Full integration checks in an environment as close as possible to the production environment. \\
\hline
\hseFill{acceptance stand} & \hseFill{composition of the acceptance stand} & Acceptance checks of the operational characteristics of the deployed system. \\
\hline
\end{tabularx}
\end{center}
\end{hseExample}

\subsection{Software used during testing}

% Подсказка: перечислите программные средства и инструменты
% (среда выполнения, фреймворки тестирования, средства контроля
% качества).
\begin{hseExample}
The following software is used during testing:

\begin{itemize}[leftmargin=1.25cm]
    \item runtime environment: \hseFill{languages and runtime environments};
    \item testing frameworks: \hseFill{testing frameworks};
    \item static analysis tools: \hseFill{linters and analysers};
    \item quality control and coverage tools: \hseFill{quality control tools};
    \item infrastructure tools: \hseFill{DBMS, brokers, containerization}.
\end{itemize}

The main way of accepting functional requirements in this work is the passing of automated tests: input data, actions and expected results are fixed in executable test scenarios, and the results of their execution are stored in \hseFill{CI/CD or test reports}.
\end{hseExample}

\subsection{Acceptance criterion for functional requirements}

\begin{hseExample}
A functional requirement with an identifier of the form \texttt{ТЗ-Ф-XX} is considered accepted if all the following conditions are met simultaneously:

\begin{enumerate}[label=\arabic*), leftmargin=1.25cm]
    \item the requirement is present in the requirements table of the Statement of Work;
    \item an automated test, test suite or end-to-end scenario is specified for the requirement;
    \item the corresponding check on the version of the program under test has completed successfully;
    \item supporting artifacts are available for the check: an execution log, a test report or a coverage report, if provided for by the given environment.
\end{enumerate}

The acceptance committee confirms the fulfilment of functional requirements in one of two ways:

% Подсказка: уточните способы подтверждения под ваш проект.
\begin{enumerate}[label=\arabic*), leftmargin=1.25cm]
    \item request access to \hseFill{project repository}, open the check results of the version under test and verify their successful completion;
    \item run the tests locally on the agreed version of the source code using the command \hseFill{test run command}.
\end{enumerate}

A local run is used as an alternative confirmation method when there is no access to CI/CD or when independent reproduction of the result is required. In case of discrepancies between the local run and CI/CD, the acceptance result is the run on the agreed version in an isolated CI/CD environment, since it fixes the code version, configuration, environment and execution artifacts.
\end{hseExample}

\subsection{Testing procedure}

% Подсказка: скорректируйте последовательность этапов под ваш проект.
\begin{hseExample}
Testing is carried out in the following order:

\begin{enumerate}[label=\arabic*), leftmargin=1.25cm]
    \item \textbf{Preparatory stage}: visual verification of the composition of the program documentation, preparation of the test environment according to sections 5.1--5.2 and deployment of the program according to the programmer's manual.
    \item \textbf{Functional acceptance}: verification of the Statement of Work requirements based on the results of the checks according to section~\ref{sec:methods}.
    \item \textbf{Static quality control}: execution and analysis of the results of \hseFill{linters and quality checks}.
    \item \textbf{Module- and service-level testing}: execution of \hseFill{testing commands or jobs} and analysis of test and coverage reports.
    \item \textbf{End-to-end testing}: execution of end-to-end scenarios in \hseFill{end-to-end testing environment}.
    \item \textbf{Recording of results}: comparison of the obtained results with the expected ones and preparation of the conclusion on the passing of the tests.
\end{enumerate}
\end{hseExample}

\subsection{Quantitative and qualitative characteristics to be assessed}

% Подсказка: скорректируйте перечень оцениваемых характеристик
% под ваш проект (наблюдаемость, нагрузка и т.д.).
\begin{hseExample}
During testing, the following characteristics are assessed:

\begin{itemize}[leftmargin=1.25cm]
    \item the percentage of successfully passed tests and checks;
    \item the absence of unforeseen failures, crashes and inconsistent data states;
    \item the correctness of \hseFill{key verifiable aspects of the program};
    \item the presence and quality of supporting test artifacts;
    \item \hseFill{additional characteristics from the Statement of Work}.
\end{itemize}
\end{hseExample}

\subsection{Test conditions}

% Подсказка: укажите необходимые условия (доступы, версии исходного
% кода, окружение, секреты).
\begin{hseExample}
Testing must be carried out subject to the following conditions:

\begin{itemize}[leftmargin=1.25cm]
    \item access to \hseFill{project repository} and check artifacts;
    \item a known identifier of the version under test: a commit identifier, a version tag or a link to a merge request;
    \item availability of \hseFill{the program runtime environment} and network access to the required resources;
    \item prepared environment variables, secrets and keys of the test environments;
    \item running the tests on the agreed version of the source code and the test configuration.
\end{itemize}
\end{hseExample}

\clearpage

% ============================================================
% 6. TEST METHODS
% ============================================================
\section{TEST METHODS}
\label{sec:methods}

\begin{hseExample}
Testing is performed by a combination of visual documentation control, automated test suites, quality control checks and acceptance checks. For each functional requirement, traceability is used down to a specific test, test suite or check from the test methods table.
\end{hseExample}

% Подсказка: для каждого требования опишите метод проверки, ожидаемый
% и фактический результаты. Заполните таблицу методов испытаний.
% Таблица методов испытаний остаётся вне жёлтого блока-образца:
% longtable внутри hseExample недопустим.

\begin{longtable}{|p{0.18\textwidth}|p{0.22\textwidth}|p{0.25\textwidth}|p{0.22\textwidth}|}
\caption{Test methods}
\label{tab:test-methods}\\
\hline
\textbf{Requirement under test} &
\textbf{Verification method} &
\textbf{Expected result} &
\textbf{Actual result} \\
\hline
\endfirsthead
\hline
\textbf{Requirement under test} &
\textbf{Verification method} &
\textbf{Expected result} &
\textbf{Actual result} \\
\hline
\endhead
\hline
\endfoot
\hline
\endlastfoot
\reqref{ТЗ-Ф-01} &
\hseFill{verification method} &
\hseFill{expected result} &
\hseFill{actual result} \\
\hline
\end{longtable}

\subsection{Testing the fulfilment of requirements for the program documentation}

Verification is performed by the visual method. For each document from clause~4.1, the following are checked:

\begin{itemize}[leftmargin=1.25cm]
    \item the presence of the document in the set;
    \item conformance to the designation, title page and structure of sections;
    \item the absence of references to outdated components, interfaces and test environments;
    \item the consistency of this document with the Statement of Work, the source listing and the programmer's manual.
\end{itemize}

\subsection{Testing of functional requirements}

% Подсказка: скорректируйте шаги под ваш CI/CD или способ запуска тестов.
\begin{hseExample}
Verification is performed by the automated method according to the requirements-to-tests traceability table. The tester performs the following actions:

\begin{enumerate}[label=\arabic*), leftmargin=1.25cm]
    \item open \hseFill{the project repository or CI/CD};
    \item select the check corresponding to the commit or version tag under test;
    \item verify that all testing jobs have completed successfully;
    \item for each requirement \texttt{ТЗ-Ф-XX}, match the specified test or test suite with a successfully completed check;
    \item save or attach screenshots and execution reports to the test materials.
\end{enumerate}

The test is considered successful if all functional requirements are traceable to an automated check, and the corresponding checks are completed without unresolved test failures. If an individual test was re-run due to environment instability, the final successful run and a link to the re-run log are recorded in the test materials.
\end{hseExample}

\subsection{Testing of static quality control and security checks}

% Подсказка: перечислите используемые линтеры, анализаторы качества
% и проверки безопасности; удалите подраздел, если они не применяются.
\begin{hseExample}
Verification is performed by the automated method. The following groups of checks are used:

\begin{itemize}[leftmargin=1.25cm]
    \item \hseFill{linters and formatters} for the analysis of formatting, linting and basic quality rules;
    \item \hseFill{quality analysis platform} for building quality analytics and aggregating test coverage;
    \item \hseFill{security scanning tools} for detecting dependency vulnerabilities and secrets in the source code.
\end{itemize}

The result of the test is the presence of the generated reports and the absence of critical blockers that make acceptance of the version impossible.
\end{hseExample}

\subsection{Testing of timing characteristics by the load method}

% Подсказка: подраздел нужен, если в ТЗ заданы временные характеристики;
% иначе удалите его.
\begin{hseExample}
Verification confirms the timing characteristics declared in the Statement of Work. The testing tool~--- \hseFill{load testing tool}.

The testing procedure:
\begin{enumerate}[label=\arabic*), leftmargin=1.25cm]
    \item prepare the workstation and deploy the program under test in the test environment;
    \item perform a load run using the command \hseFill{run start command};
    \item perform a stepped series of runs with the parameters \hseFill{number of users and duration}; treat the first \hseFill{warm-up period duration} as the warm-up period and exclude it from the acceptance metrics;
    \item record the summary indicators of the run: the total number of requests, the number of failures, the aggregated request rate, as well as the response time percentiles (p50, p95, p99);
    \item save or attach to the test materials the report and screenshots of the summary statistics.
\end{enumerate}

Success criteria:
\begin{itemize}[nosep, leftmargin=1.25cm]
    \item the aggregated request rate after the warm-up period is at least \hseFill{target request rate};
    \item the proportion of successful responses after the warm-up period is not lower than \hseFill{target proportion of successful responses};
    \item there are no unhandled exceptions and service crashes during the run;
    \item a report is generated with a table of requests, response time statistics and load graphs.
\end{itemize}
\end{hseExample}

\subsection{Testing of requirements traceability and artifact completeness}

\begin{hseExample}
For each functional requirement of the Statement of Work with an identifier of the form \texttt{ТЗ-Ф-XX}, an automated check and the location of the corresponding test must be specified. Non-functional requirements are confirmed by separate test methods: \hseFill{verification methods for non-functional requirements}. Additionally, the completeness of the supporting test materials attached to this document is verified.
\end{hseExample}

\clearpage
% Страницы источников печатаются без нижнего штампа (как в реальных
% комплектах) — только номер листа и код документа сверху.
\pagestyle{appendix}

% ============================================================
% REFERENCES (IEEE style)
% ============================================================
\section*{REFERENCES}
\addcontentsline{toc}{section}{REFERENCES}

% Подсказка: добавьте ПЕРЕД ГОСТами источники по технологиям вашего
% проекта (IEEE-стиль): A.~B.~Author, \enquote{Title,} venue, year.
\renewcommand{\refname}{}
\begin{thebibliography}{99}
\bibitem{gost19101}
\emph{GOST 19.101-77. Types of Programs and Program Documents.} Unified System for Program Documentation. Moscow, Russia: IPK Izdatelstvo Standartov, 2001.

\bibitem{gost19103}
\emph{GOST 19.103-77. Designation of Programs and Program Documents.} Unified System for Program Documentation. Moscow, Russia: IPK Izdatelstvo Standartov, 2001.

\bibitem{gost19104}
\emph{GOST 19.104-78. Basic Inscriptions.} Unified System for Program Documentation. Moscow, Russia: IPK Izdatelstvo Standartov, 2001.

\bibitem{gost19105}
\emph{GOST 19.105-78. General Requirements for Program Documents.} Unified System for Program Documentation. Moscow, Russia: IPK Izdatelstvo Standartov, 2001.

\bibitem{gost19106}
\emph{GOST 19.106-78. Requirements for Program Documents Produced in Printed Form.} Unified System for Program Documentation. Moscow, Russia: IPK Izdatelstvo Standartov, 2001.

\bibitem{gost19301}
\emph{GOST 19.301-79. Test Program and Methodology. Requirements for Content and Presentation.} Unified System for Program Documentation. Moscow, Russia: IPK Izdatelstvo Standartov, 2001.

\bibitem{gost19604}
\emph{GOST 19.604-78. Rules for Making Changes to Program Documents Produced in Printed Form.} Unified System for Program Documentation. Moscow, Russia: IPK Izdatelstvo Standartov, 2001.
\end{thebibliography}

\clearpage


% ============================================================
% ЛИСТ РЕГИСТРАЦИИ ИЗМЕНЕНИЙ
% ============================================================
\input{../common/change_log}

\end{document}
